% -*- coding: utf-8 -*-
% ===================================================================
% DODATEK T6: Perturbacyjne rozwiniecie eta(delta) do O(delta^2):
%             alpha_3 = pi^2/128 + P_cos
% ===================================================================
% Math operator shortcuts (local to this appendix)
\providecommand{\Ci}{\operatorname{Ci}}
\providecommand{\Si}{\operatorname{Si}}
\providecommand{\Li}{\operatorname{Li}}

\section*{Dodatek T6: drugi wsp\'o\l{}czynnik rozwini\k{e}cia $\eta(\delta)$: $\alpha_3 \approx 0.08972$}%
\addcontentsline{toc}{section}{Dodatek T6: $\alpha_3 = \pi^2/128 + P_{\cos}$ (sub-breakthrough)}%
\label{app:T6}

\subsection*{Streszczenie}

W~dodatku~\ref{app:T5} wyprowadzili\'smy analityczny wz\'or dla sta\l{}ej
asymetrii deficit/excess: $c_1 = 1 - \ln(3)/4$. W~niniejszym dodatku
pokazujemy, \.ze rozwini\k{e}cie funkcji $\eta(\delta) = A_{\rm tail}(1+\delta)/|\delta|$
w~okolicy $\delta = 0$ ma rygorystyczn\k{a} struktur\k{e}:
\begin{equation}
  \boxed{\;\eta(\delta) = 1 + \tfrac{c_1}{2}\,\delta + \alpha_3\,\delta^2 + O(\delta^3)\;}
  \label{eq:eta-expansion-full}
\end{equation}
gdzie wsp\'o\l{}czynnik kwadratowy $\alpha_3$ pochodzi z~rozk\l{}adu
perturbacyjnego w~O($\delta^2$) i~O($\delta^3$) i~rozpada si\k{e} kanonicznie na:
\begin{equation}
  \alpha_3 = \underbrace{\tfrac{I_{\sin}^2}{2}}_{\pi^2/128\; \text{(\textbf{udow.})}}
           + \underbrace{P_{\cos}}_{0.012615939\ldots\;\text{(\textbf{numerycznie})}}.
  \label{eq:alpha3-decomp}
\end{equation}
Cz\l{}on $\pi^2/128$ zosta\l{} udowodniony analitycznie. Sta\l{}a $P_{\cos}$
jest obliczona numerycznie z~precyzj\k{a} 30 cyfr znacz\k{a}cych
(test zbie\.zno\'sci dps=25/40/60 --- stabilna), lecz nie posiada
zidentyfikowanej formy zamkni\k{e}tej.

\medskip
\textbf{Rewizja:} Wcze\'sniejsza hipoteza $\alpha_3 = \pi^2/110$ (oparta na
ekstrapolacji Richardsona z~precyzj\k{a} $\sim 10^{-5}$) zosta\l{}a \textbf{obalona}
przez niezale\.zn\k{a} weryfikacj\k{e} sw\'apow\k{a} przy dps=60. R\'o\.znica wynosi
$\alpha_3 - \pi^2/110 = -1.4527 \times 10^{-6}$ i~jest \emph{stabilna}
we~wszystkich precyzjach. Wyszukiwanie PSLQ z~baz\k{a}
$\{\pi^2, \pi^4, \ln 2, \ln 3, \ln^2 2, \ln^2 3, \ln 2\cdot\ln 3,
\zeta(3), G, \chi_2(1/3), \chi_2(1/2), 1\}$ nie~znajduje relacji
o~wsp\'o\l{}czynnikach ca\l{}kowitych do $10^{10}$.

% -------------------------------------------------------------------
\subsection{Struktura rozwini\k{e}cia $\eta(\delta)$}
\label{sec:T6-structure}
% -------------------------------------------------------------------

Niech $g_0 = 1 + \delta$ i~$g(r; \delta)$ spe\l{}nia substratowe ODE TGP
\eqref{eq:substrate-ODE}. Definiujemy
\begin{equation}
  \eta(\delta) := \frac{A_{\rm tail}(1+\delta)}{|\delta|},
\end{equation}
gdzie $A_{\rm tail}$ jest amplitud\k{a} asymptotyk\k{a} ogona,
$(g-1) \sim A_{\rm tail}\sin(r+\varphi)/r$.

Z~symetrii perturbacji (zamiana $\delta \leftrightarrow -\delta$
odpowiada zamianie excess $\leftrightarrow$ deficit), zachodzi:
\begin{align}
  \eta_{\rm exc}(\delta) &= 1 + \alpha_2\,\delta + \alpha_3\,\delta^2 + \alpha_4\,\delta^3 + \alpha_5\,\delta^4 + \dots, \\
  \eta_{\rm def}(\delta) &= 1 - \alpha_2\,\delta + \alpha_3\,\delta^2 - \alpha_4\,\delta^3 + \alpha_5\,\delta^4 - \dots.
\end{align}
St\k{a}d sumatywny (symetryczny) sk\l{}adnik:
\begin{equation}
  \eta_{\rm sym}(\delta) := \tfrac{1}{2}(\eta_{\rm exc} + \eta_{\rm def})
    = 1 + \alpha_3\,\delta^2 + \alpha_5\,\delta^4 + O(\delta^6).
\end{equation}
Znane wyniki:
\begin{itemize}
  \item $\alpha_2 = c_1/2 = \tfrac{1}{2}(1 - \ln 3/4)$ --- \textbf{udowodnione}
        w~dodatku~\ref{app:T5} (Frullani).
  \item $\alpha_3 = \pi^2/128 + P_{\cos} = 0.089722223674\ldots$ --- \textbf{obliczone
        numerycznie do 30 cyfr}; brak zidentyfikowanej formy zamkni\k{e}tej
        (niniejszy dodatek).
\end{itemize}

% -------------------------------------------------------------------
\subsection{Perturbacja O($\delta^2$) + O($\delta^3$)}
\label{sec:T6-perturbation}
% -------------------------------------------------------------------

Rozwijamy $g(r; \delta) = 1 + \delta\,f(r) + \delta^2\,h(r) + \delta^3\,p(r) + O(\delta^4)$.

\subsection*{O($\delta$) --- znane z~dodatku T5}

$f(r) = \sin(r)/r = j_0(r)$ spe\l{}nia
\begin{equation}
  f'' + \tfrac{2}{r}f' + f = 0,
\end{equation}
z~asymptotyk\k{a} $f(r) \sim \sin(r)/r$, wi\k{e}c $rf \sim \sin r$ i~wk\l{}ad do
amplitudy ogona: $\delta\cdot\sin(r)$.

\subsection*{O($\delta^2$) --- r\'ownanie dla $h$}

Rozwijaj\k{a}c $(1/g)(g')^2$ i~$-g$ do drugiego rz\k{e}du:
\begin{equation}
  h'' + \tfrac{2}{r}h' + h = -(f')^2.
  \label{eq:T6-ODE-h}
\end{equation}
Substytuuj\k{a}c $u(r) = r\,h(r)$, dostajemy prostsz\k{a} form\k{e}:
\begin{equation}
  u'' + u = -r\,(f'(r))^2.
  \label{eq:T6-u-eq}
\end{equation}
Wariacja parametr\'ow z~fundamentalnymi rozwi\k{a}zaniami $\sin(r), -\cos(r)$
daje asymptotycznie
\begin{equation}
  u(r) \xrightarrow{r \to \infty} \sin(r)\,I_{\cos} - \cos(r)\,I_{\sin},
\end{equation}
gdzie (por. dodatek~T5):
\begin{align}
  I_{\cos} &= \int_0^{\infty}\cos(s)\cdot(-s\,(f'(s))^2)\,ds = \tfrac{1}{2} - \tfrac{\ln 3}{8}, \\
  I_{\sin} &= \int_0^{\infty}\sin(s)\cdot(-s\,(f'(s))^2)\,ds = -\tfrac{\pi}{8}.
\end{align}

\subsection*{O($\delta^3$) --- r\'ownanie dla $p$}

Rozwini\k{e}cie ODE do trzeciego rz\k{e}du daje ni-liniowe sprz\k{e}\.zenie
przez \textbf{krzy\.zowe cz\l{}ony} $f\,h$:
\begin{equation}
  p'' + \tfrac{2}{r}p' + p = -2 f'(r)\,h'(r) + f(r)\,(f'(r))^2.
  \label{eq:T6-ODE-p}
\end{equation}
Substytucja $v = r\,p$ upraszcza do:
\begin{equation}
  v'' + v = r\big[-2 f'(r)\,h'(r) + f(r)\,(f'(r))^2\big] =: r\,R(r).
\end{equation}
Asymptotycznie $v(r) \to \sin(r)\,P_{\cos} - \cos(r)\,P_{\sin}$ z:
\begin{equation}
  P_{\cos} = \int_0^{\infty} \cos(s)\cdot s\,R(s)\,ds.
  \label{eq:T6-Pcos-def}
\end{equation}

\subsection*{Rozk\l{}ad $\alpha_3$}

Z~rozwini\k{e}cia $A_{\rm tail} = \sqrt{S^2 + C^2}$, gdzie
\begin{align}
  S &= \delta + \delta^2\,I_{\cos} + \delta^3\,P_{\cos} + \dots, \\
  C &= -\delta^2\,I_{\sin} - \delta^3\,P_{\sin} + \dots,
\end{align}
obliczamy (dla $\delta > 0$):
\begin{equation}
  \eta(\delta) = \frac{A}{|\delta|} = 1 + I_{\cos}\,\delta + \left(\tfrac{I_{\sin}^2}{2} + P_{\cos}\right)\delta^2 + O(\delta^3).
  \label{eq:T6-alpha3-full}
\end{equation}
St\k{a}d:
\begin{equation}
  \boxed{\;\alpha_3 = \tfrac{I_{\sin}^2}{2} + P_{\cos}.\;}
\end{equation}
Podstawiaj\k{a}c $I_{\sin} = -\pi/8$:
\begin{equation}
  \tfrac{I_{\sin}^2}{2} = \tfrac{\pi^2/64}{2} = \tfrac{\pi^2}{128}.
\end{equation}

% -------------------------------------------------------------------
\subsection{Weryfikacja numeryczna}
\label{sec:T6-numerics}
% -------------------------------------------------------------------

Dwie niezale\.zne procedury da\l{}y konsystentny wynik.

\subsection*{Metoda 1: pomiar $\alpha_3$ z~pe\l{}nego ODE}

Skrypt \texttt{r6\_c2\_ultra\_precision.py} rozwi\k{a}zuje pe\l{}ne ODE dla
szeregu $r_{\max} \in \{400, 600, 800, 1200, 1600\}$ i~stosuje ekstrapolacj\k{e}
Richardsona $1/r_{\max}^2 \to 0$:
\begin{center}
\begin{tabular}{r|rr}
\hline
$r_{\max}$ & $\alpha_3$ & diff od $\pi^2/110$ \\
\hline
 400  & 0.089710060 & $-1.4 \times 10^{-5}$ \\
 600  & 0.089714669 & $-9.0 \times 10^{-6}$ \\
 800  & 0.089717060 & $-6.6 \times 10^{-6}$ \\
1200  & 0.089722723 & $-9.5 \times 10^{-7}$ \\
1600  & 0.089722417 & $-1.3 \times 10^{-6}$ \\
\hline
$r_{\max} \to \infty$ & \textbf{0.089722365} & $+1.3 \times 10^{-6}$ \\
teoria $\pi^2/110$ & 0.089723676 & --- \\
\hline
\end{tabular}
\end{center}

Walidacja metody: analogiczna ekstrapolacja dla $\alpha_2$ daje
$\alpha_2 = 0.362666$ vs~$c_1/2 = 0.362673$ (diff $7 \times 10^{-6}$).

\subsection*{Metoda 2: niezale\.zna perturbacja $h, p$}

Skrypt \texttt{r6\_c3\_perturbative\_h\_and\_p.py} rozwi\k{a}zuje
\eqref{eq:T6-u-eq} dla $u(r)$, interpoluje, a~nast\k{e}pnie rozwi\k{a}zuje
\eqref{eq:T6-ODE-p} dla $v(r) = r p(r)$.

\begin{center}
\begin{tabular}{lrrr}
\hline
Wielko\'s\'c & Numeryczna (scipy) & Teoria/Ref. & diff \\
\hline
$I_{\cos}$ (O($\delta$))  & $+0.362674575$ & $\frac{1}{2} - \frac{\ln 3}{8}$ (udow.) & $+1.1 \times 10^{-6}$ \\
$I_{\sin}$ (O($\delta$))  & $-0.392698433$ & $-\pi/8$ (udow.)                        & $+6.5 \times 10^{-7}$ \\
$P_{\cos}$ (O($\delta^3$)) & $+0.012617724$ & \textbf{0.012615939} (mpmath 30 cyfr)   & $+1.8 \times 10^{-6}$ \\
$I_{\sin}^2/2 + P_{\cos}$ & $+0.089724009$ & \textbf{0.089722224} (mpmath 30 cyfr)   & $+1.8 \times 10^{-6}$ \\
\hline
\end{tabular}

\medskip
\textbf{Uwaga}: r\'o\.znica $1.8 \times 10^{-6}$ mi\k{e}dzy scipy (float64,
$r_{\max}=1800$) a~mpmath (dps=60, quadosc) jest artefaktem precyzji scipy.
Warto\'s\'c referencyjn\k{a} \textbf{30-cyfrowa} pochodzi z~niezale\.znego testu
zbie\.zno\'sci (dps=25/40/60 --- identyczna do 30 cyfr):
\begin{equation}
  \alpha_3 = 0.0897222236736253260474946788048\ldots
\end{equation}
\end{center}

\subsection*{Podsumowanie dekompozycji $\alpha_3$}

\begin{equation}
  \alpha_3 = \tfrac{I_{\sin}^2}{2} + P_{\cos} = \underbrace{\tfrac{\pi^2}{128}}_{\text{udow.}}
   + \underbrace{P_{\cos}}_{= K_c^{(I)} - 2K_c^{(II)}}.
\end{equation}
Warto\'sci numeryczne (30 cyfr, stabilne dla dps=25/40/60):
\begin{align}
  P_{\cos} &= 0.012615939290114711837850\ldots, \notag \\
  K_c^{(II)} &= -0.031879037931728580134156\ldots, \notag \\
  \alpha_3 &= 0.089722223673625326047495\ldots \notag
\end{align}
Wcze\'sniejsza hipoteza $\alpha_3 = \pi^2/110 \approx 0.089723676$ jest \textbf{obalona}
(r\'o\.znica $-1.453 \times 10^{-6}$, stabilna).

% -------------------------------------------------------------------
\subsection{Implikacje fizyczne}
\label{sec:T6-physics}
% -------------------------------------------------------------------

\subsection*{Pe\l{}ny wz\'or dla $\eta(\delta)$}

\begin{equation}
  \boxed{\;\eta(\delta) = 1 + \tfrac{1-\ln 3/4}{2}\,\delta + \alpha_3\,\delta^2 + O(\delta^3)\;}
\end{equation}
gdzie $\alpha_3 = \pi^2/128 + P_{\cos} = 0.089722224\ldots$ (30 cyfr znacz\k{a}cych).
Pierwsza sta\l{}a ($c_1/2 = (1 - \ln 3/4)/2$) udowodniona analitycznie;
druga ($\alpha_3$) obliczona z~wysok\k{a} precyzj\k{a}, lecz bez zamkni\k{e}tej formy.

\subsection*{Korekta boost masy}

Dla elektronu ($\delta_e = -0.131$) i~mionu ($\delta_\mu = +0.407$):
\begin{align*}
  \eta_e &= 1 + c_1/2 \cdot (-0.131) + \alpha_3 \cdot 0.01716 + \ldots \approx 0.9537, \\
  \eta_\mu &= 1 + c_1/2 \cdot (+0.407) + \alpha_3 \cdot 0.1656 + \ldots \approx 1.1627.
\end{align*}
Stosunek $\eta_\mu/\eta_e \approx 1.219$, co przyczynia si\k{e} do boost
$(\eta_\mu/\eta_e)^4 \approx 2.21$ (target PDG: $2.195$).

\subsection*{Struktura $P_{\cos}$}

Sta\l{}a $P_{\cos}$ jest definiowana przez podw\'ojn\k{a} ca\l{}k\k{e}:
\begin{equation}
  P_{\cos} = \int_0^{\infty} \cos(s)\,s\,\big[-2 f'(s)\,h'(s) + f(s)\,(f'(s))^2\big]\,ds
          = 0.012615939290\ldots
\end{equation}
i~rozk\l{}ada si\k{e} na $P_{\cos} = K_c^{(I)} - 2K_c^{(II)}$, gdzie
$K_c^{(I)} = (\ln 2 - 1)/6$ (udowodnione) a~$K_c^{(II)} = -0.031879038\ldots$
Poszukiwanie PSLQ w~bazie 15 sta\l{}ych matematycznych
(w\l{}\k{a}cznie z~$\pi^2, \pi^4, \zeta(3), G, \chi_2(1/3)$)
przy wsp\'o\l{}czynnikach do $10^{10}$ \textbf{nie~znalaz\l{}o \.zadnej relacji}
--- $P_{\cos}$ jest prawdopodobnie sta\l{}\k{a} ``now\k{a}'', specyficzn\k{a}
dla substratowego ODE TGP.

% -------------------------------------------------------------------
\subsection{Post\k{e}p analityczny: dow\'od cz\k{e}\'sciowy $P_{\cos}$ przez IBP}
\label{sec:T6-IBP}
% -------------------------------------------------------------------

W~niniejszej sekcji prezentujemy nowe wyniki analityczne (sesja 2026-04-16)
daj\k{a}ce dwie nowe to\.zsamo\'sci domkni\k{e}te oraz cz\k{e}\'sciow\k{a} dekompozycj\k{e}
$P_{\cos}$.

\subsubsection*{To\.zsamo\'s\'c IBP z~r\'ownania Bessela sferycznego}

Kluczow\k{a} to\.zsamo\'sci\k{a} derivowan\k{a} z~$f'' + (2/s)f' + f = 0$ ($f = \sin(s)/s$) jest:
\begin{equation}
  \boxed{\;s\,(f')^2 = \frac{d}{ds}\big[s\,f\,f'\big] + s\,f^2 + \tfrac{1}{2}\frac{d}{ds}\big[f^2\big].\;}
  \label{eq:T6-IBP-identity}
\end{equation}
Sprawdzenie: $\frac{d}{ds}[s f f'] = f f' + s (f')^2 + s f f''$;
podstawiaj\k{a}c $s f f'' = -s f^2 - 2 f f'$ (z~ODE), dostajemy
$\frac{d}{ds}[s f f'] = s (f')^2 - s f^2 - f f'$, sk\k{a}d~\eqref{eq:T6-IBP-identity}
przez $f f' = (1/2)(f^2)'$.

\subsubsection*{Czysty dow\'od $J_c$ i~$J_s$}

Niech $J_c = \int_0^\infty \cos(s)\,s\,(f')^2\,ds$ i~$J_s = \int_0^\infty \sin(s)\,s\,(f')^2\,ds$
(zwi\k{a}zek: $\alpha_2 = -J_c$, $I_{\sin}^2/2 = J_s^2/2$). Stosuj\k{a}c~\eqref{eq:T6-IBP-identity}
oraz IBP z~$\cos(s)$:
\begin{equation}
  J_c = \int_0^\infty \sin(s)\cdot s\,f\,f'\,ds + \int_0^\infty \cos(s)\cdot s\,f^2\,ds
        - \tfrac{1}{2} + \tfrac{1}{2}\int_0^\infty \sin(s)\cdot f^2\,ds.
\end{equation}
(wk\l{}ad $-1/2$ pochodzi z~$\cos(0)\,f^2(0)/2 = 1/2$ w~granicy boundary.)

Trzy pozosta\l{}e ca\l{}ki redukuj\k{a} si\k{e} do Frullani/Dirichleta:
\begin{align}
  \int_0^\infty \sin(s)\,s f f'\,ds &= \int_0^\infty \sin^2(s)\,f'(s)\,ds
     = -\tfrac{\ln 3}{2}, \\
  \int_0^\infty \cos(s)\,s f^2\,ds &= \int_0^\infty \tfrac{\cos(s)\sin^2 s}{s}\,ds
     = \tfrac{\ln 3}{4}, \\
  \int_0^\infty \sin(s)\,f^2\,ds &= \int_0^\infty \tfrac{\sin^3 s}{s^2}\,ds
     = \tfrac{3\ln 3}{4}.
\end{align}
Pierwsze dwie z~klasycznej to\.zsamo\'sci $\int [cos(s) - cos(3s)]/s\,ds = \ln 3$;
trzecia z~IBP do $3\int \sin^2(s)\cos(s)/s\,ds = (3/4)\ln 3$.

Sumuj\k{a}c:
\begin{equation}
  \boxed{\;J_c = -\tfrac{\ln 3}{2} + \tfrac{\ln 3}{4} - \tfrac{1}{2} + \tfrac{3\ln 3}{8}
          = \tfrac{\ln 3}{8} - \tfrac{1}{2}.\;}
\end{equation}
Analogicznie (z~IBP dla $\sin$, wyliczenia Dirichleta, $I(3) = \pi$ dla $\int [\cos s - \cos 3s]/s^2 ds$):
\begin{equation}
  \boxed{\;J_s = 0 + \tfrac{\pi}{4} - \tfrac{1}{2}\cdot\tfrac{\pi}{4} = \tfrac{\pi}{8}.\;}
\end{equation}
St\k{a}d \textbf{czysty dow\'od} $\alpha_2 = -J_c = 1/2 - \ln(3)/8 = I_{\cos}$ oraz
$I_{\sin}^2/2 = J_s^2/2 = \pi^2/128$.

Weryfikacja numeryczna (mpmath dps=40, \texttt{r6\_c4\_proof\_Jc\_Js.py}):
wszystkie 6 ca\l{}ek elementarnych zgadza si\k{e} do $10^{-42}$.

\subsubsection*{Dekompozycja $P_{\cos}$ i~wynik $K_c^{(I)}$}

Rozbijaj\k{a}c $R(s) = f(f')^2 - 2 f' h'$ na dwa cz\l{}ony:
\begin{equation}
  P_{\cos} = K_c^{(I)} - 2 K_c^{(II)},
\end{equation}
gdzie:
\begin{align}
  K_c^{(I)} &:= \int_0^\infty \cos(s)\,s\,f(s)\,(f'(s))^2\,ds, \\
  K_c^{(II)} &:= \int_0^\infty \cos(s)\,s\,f'(s)\,h'(s)\,ds.
\end{align}

Do $K_c^{(I)}$ stosujemy to\.zsamo\'s\'c:
\begin{equation}
  s\,f\,(f')^2 = \tfrac{1}{2}\tfrac{d}{ds}\big[s\,f^2\,f'\big] + \tfrac{1}{6}\tfrac{d}{ds}\big[f^3\big] + \tfrac{1}{2}\,s\,f^3,
\end{equation}
(pochodna od tej samej ODE) i~IBP z~$\cos(s)$:
\begin{equation}
  K_c^{(I)} = \tfrac{1}{2}\int \sin(s)\,s f^2 f'\,ds - \tfrac{1}{6}
              + \tfrac{1}{6}\int \sin(s)\,f^3\,ds + \tfrac{1}{2}\int \cos(s)\,s\,f^3\,ds.
\end{equation}
Potrzebne dwie ca\l{}ki elementarne:
\begin{align}
  \int_0^\infty \tfrac{\sin^3(s)\cos(s)}{s^2}\,ds &= \tfrac{\ln 2}{2} \quad (\text{Frullani dla } \cos 2s - \cos 4s), \\
  \int_0^\infty \tfrac{\sin^4(s)}{s^3}\,ds &= \ln 2 \quad (\text{IBP i~powy\.zsza}).
\end{align}

Razem:
\begin{equation}
  \boxed{\;K_c^{(I)} = \tfrac{\ln 2 - 1}{6} \approx -0.05114.\;}
  \label{eq:T6-KcI}
\end{equation}
Zweryfikowane numerycznie do 40 cyfr (mpmath quadosc).

\subsubsection*{Wynik $K_s^{(I)} = \pi/12$ \textbf{(NOWE, udow.)}}

Analogicznie definiujemy $K_s^{(I)} := \int_0^\infty \sin(s)\,s\,f(s)\,(f'(s))^2\,ds$
(potrzebne do $P_{\sin}$, sk\l{}adowej fazy ogona O($\delta^3$) --- patrz
sekcja~\ref{sec:T6-alpha5}).

Rozwijaj\k{a}c $s\,f\,(f')^2 = \sin(s)(s\cos(s)-\sin(s))^2/s^4$:
\begin{equation}
  K_s^{(I)} = \underbrace{\int_0^\infty \frac{\sin^2(s)\cos^2(s)}{s^2}\,ds}_{J_1}
            - 2\underbrace{\int_0^\infty \frac{\sin^3(s)\cos(s)}{s^3}\,ds}_{J_2}
            + \underbrace{\int_0^\infty \frac{\sin^4(s)}{s^4}\,ds}_{J_3}.
\end{equation}
Ka\.zda ca\l{}ka jest elementarna:
\begin{itemize}
  \item $J_1$: Z~$\sin^2\cos^2 = \sin^2(2s)/4$ i~$\int_0^\infty \sin^2(as)/s^2\,ds = \pi a/2$
        wynika $J_1 = \tfrac{1}{4}\cdot \pi\cdot 2/1 \cdot 1/2 = \pi/4$.
  \item $J_2$: Z~$\sin^3\cos = (\sin 2s - \sin 4s/2)/4$
        i~wzoru $\int_0^\infty (\sin(as)-as)/s^3\,ds = -\pi a^2/4$
        (wyprowadzonego przez $I'(a) = -\int (1-\cos as)/s^2\,ds = -\pi|a|/2$)
        dostajemy $J_2 = \pi/4$.
  \item $J_3 = \int_0^\infty \operatorname{sinc}^4(s)\,ds = \pi/3$ --- wynik klasyczny.
\end{itemize}
Razem:
\begin{equation}
  \boxed{\;K_s^{(I)} = \tfrac{\pi}{4} - 2\cdot\tfrac{\pi}{4} + \tfrac{\pi}{3}
          = -\tfrac{\pi}{4} + \tfrac{\pi}{3} = \tfrac{\pi}{12}.\;}
  \label{eq:T6-KsI}
\end{equation}
Zweryfikowane numerycznie do 32 cyfr.

\subsubsection*{Pozostaj\k{a}cy krok: $K_c^{(II)}$}

Wz\'or \eqref{eq:T6-KcI} daje:
\begin{equation}
  P_{\cos} = K_c^{(I)} - 2 K_c^{(II)} = \tfrac{\ln 2 - 1}{6} - 2 K_c^{(II)}.
\end{equation}
Numerycznie $K_c^{(II)} = -0.031879037931728580\ldots$ (30 cyfr), co jest
potwierdzone przez niezale\.zn\k{a} \emph{sw\'apow\k{a}} dekompozycj\k{e}
(rozdzia\l{}~\ref{sec:T6-swap}), stabiln\k{a} przy dps=25/40/60.

% -------------------------------------------------------------------
\subsubsection{Sw\'ap kolejno\'sci podwójnej ca\l{}ki: $K_c^{(II)} = -A_1 - A_2 + A_3 - A_4$}
\label{sec:T6-swap}
% -------------------------------------------------------------------

Z $h$-r\'ownania $h'' + (2/r) h' + h = -(f')^2$, przez funkcj\k{e} Greena uzyskujemy
$u(r) = \sin(r) J_c(r) - \cos(r) J_s(r)$ (gdzie $u \equiv -rh$),
$u'(r) = \cos(r) J_c(r) + \sin(r) J_s(r)$, a~st\k{a}d
$h'(r) = -u'(r)/r + u(r)/r^2$. Podstawiaj\k{a}c do definicji $K_c^{(II)}$
i~zamieniaj\k{a}c kolejno\'s\'c ca\l{}kowania (Fubini):
\begin{equation}
  K_c^{(II)} = -A_1 - A_2 + A_3 - A_4,
  \qquad A_i = \int_0^\infty \tau_i(t)\, t (f'(t))^2\, \Phi_i(t)\, dt,
\end{equation}
gdzie $\tau_1 = \tau_3 = \cos$, $\tau_2 = \tau_4 = \sin$, a~\emph{tail functions}
$\Phi_i(t) = \int_t^\infty \psi_i(s)\, ds$ dane s\k{a} przez kombinacje
$Ci, Si, \sin, \cos$ o~argumentach $t$ i~$3t$:
\begin{align}
  \Phi_1(t) &= -\tfrac{1}{2} \Ci(t) + \tfrac{1}{2} \Ci(3t) - \tfrac{\sin(t)+\sin(3t)}{4t}, \\
  \Phi_2(t) &= \tfrac{1}{2}\bigl[\Si(3t) - \Si(t)\bigr] - \tfrac{\cos(t) - \cos(3t)}{4t}, \\
  \Phi_3(t) &= \tfrac{3\sin(t) - \sin(3t)}{8t} + \tfrac{3}{8}[\Ci(3t) - \Ci(t)]
              + \tfrac{\cos(3t) - \cos(t)}{8 t^2}, \\
  \Phi_4(t) &= \tfrac{5\cos(t) - \cos(3t)}{8t} - \tfrac{\pi}{8} + \tfrac{5}{8}\Si(t)
              - \tfrac{3}{8}\Si(3t) - \tfrac{\sin(t) + \sin(3t)}{8 t^2}.
\end{align}
Powy\.zsze uzyskane s\k{a} przez elementarne IBP~dla funkcji $\cos^2, \cos\sin$
pomno\.zonych przez $f'(s)/s^k$ ($k \in \{0,1\}$). Wszystkie cztery wzory
zweryfikowane numerycznie do \textbf{30+ cyfr} przy pi\k{e}ciu r\'o\.znych warto\'sciach $t$
(skrypt \texttt{r6\_c6\_KcII\_phi\_swap.py}).

Numeryczna ewaluacja sw\'apowej dekompozycji (dps=50) daje:
\begin{align}
  A_1 &= 0.07031453567\ldots, &
  A_2 &= 0.02918815038\ldots, \\
  A_3 &= 0.00454851497\ldots, &
  A_4 &= -0.06307513316\ldots, \\
  K_c^{(II)} &= 0.031879038\ldots \text{(z minusem: $-0.031879038$)}.
\end{align}
Warto\'sci stabilne we~wszystkich testowanych precyzjach (dps=25, 40, 60 ---
identyczne do 30 cyfr). Niezale\.zna weryfikacja metod\k{a} sumowania
p\'o\l{}-okres\'ow z~akceleracj\k{a} Wynna-Epsilona potwierdza zbie\.zno\'s\'c
ku~tym samym warto\'sciom (skrypt \texttt{r6\_c16}).
Powy\.zsze wyklucza hipotez\k{e} $\alpha_3 = \pi^2/110$:
r\'o\.znica wynosi $-1.453 \times 10^{-6}$ i~\emph{nie maleje z~rosn\k{a}c\k{a} precyzj\k{a}}.

\subsubsection*{Stolik ca\l{}ek momentowych}

Do analitycznej ewaluacji $A_i$ potrzebne s\k{a}:
\begin{align}
  M_1(a,b) &:= \int_0^\infty \tfrac{\sin(a t)\, \Si(b t)}{t}\, dt
            = \begin{cases} \chi_2(b/a), & a > b \\ \pi^2/8, & a = b \\ \pi^2/4 - \chi_2(a/b), & a < b \end{cases} \\
  M_2(a,b) &:= \int_0^\infty \tfrac{\sin(a t)\, \Ci(b t)}{t}\, dt
            = \begin{cases} -\tfrac{\pi}{2} \ln(a/b), & a > b \\ 0, & a \le b \end{cases} \\
  M_3(a,b) &:= \int_0^\infty \tfrac{\cos(a t)\, \Si(b t)}{t}\, dt
            = \begin{cases} \tfrac{\pi}{2} \ln(b/a), & b > a \\ 0, & b \le a \end{cases}
\end{align}
gdzie $\chi_2(x) = x + x^3/9 + x^5/25 + \cdots = \tfrac{1}{2}[\Li_2(x) - \Li_2(-x)]$
jest funkcj\k{a} chi Legendre'a. Szczeg\'oln\k{a} warto\'sci\k{a} jest
$\chi_2(1/3) = 0.337623178\ldots$ --- ta sta\l{}a pojawi si\k{e} w~analitycznych
warto\'sciach $A_1$ i~$A_2$.

Zweryfikowane numerycznie: $M_1(1,3) + M_1(3,1) = \pi^2/4$ (symetria IBP),
$M_1(3,1) = \chi_2(1/3)$, $M_3(a,b) = -M_2(b,a)$ (dualno\'s\'c Ci$\leftrightarrow$Si).

% -------------------------------------------------------------------
\subsection{Wsp\'o\l{}czynnik czwartego rz\k{e}du: $\alpha_5 \approx 0.0275$}
\label{sec:T6-alpha5}
% -------------------------------------------------------------------

Nast\k{e}pny wsp\'o\l{}czynnik parzysty w~rozwini\k{e}ciu
$\eta_{\rm sym}(\delta) = 1 + \alpha_3 \delta^2 + \alpha_5 \delta^4 + O(\delta^6)$
wyznaczamy z~pe\l{}nego nieliniowego ODE substratowego.

\subsubsection*{Metoda}

Rozwi\k{a}zujemy ODE dla $g_0 = 1 \pm \delta$ przy wielu warto\'sciach
$\delta \in [0.04, 0.25]$ (wystarczaj\k{a}co du\.ze, by~sygna\l{} $\delta^4$
przekracza\l{} szum float64), ekstrapolujemy amplitud\k{e} $A_{\rm tail}$
z~okna $[80, r_{\max}-80]$ z~korekcj\k{a} rozk\l{}adu $1/r$, i~dopasowujemy
$(\eta_{\rm sym} - 1)/\delta^2 = \alpha_3 + \alpha_5 \delta^2 + \alpha_7 \delta^4 + \alpha_9 \delta^6$
wielomianem stopnia 3 w~$\delta^2$.

\subsubsection*{Wynik}

Pomiary przy $r_{\max}$ od 400 do 3000, z~ekstrapolacj\k{a} Richardsona:
\begin{center}
\begin{tabular}{lcc}
\hline
$r_{\max}$ & $\alpha_5$ (fit deg~3) & b\l{}\k{a}d $\alpha_3$ \\
\hline
$400$  & 0.027464 & $+2.9 \times 10^{-5}$ \\
$600$  & 0.027485 & $-5.1 \times 10^{-6}$ \\
$1000$ & 0.027495 & $-9.5 \times 10^{-6}$ \\
$2000$ & 0.027508 & $-7.2 \times 10^{-6}$ \\
$3000$ & 0.027513 & $-3.8 \times 10^{-6}$ \\
\hline
Richardson $(2000,3000)$ $1/r^2$ & \textbf{0.02752} & --- \\
\hline
\end{tabular}
\end{center}
Wniosek: $\alpha_5 \approx 0.0275 \pm 0.0001$ (4~cyfry znacz\k{a}ce).
Dodatkowe szacunki: $\alpha_7 \approx 0.011$, $\alpha_9 \approx 0.007$.

\subsubsection*{Struktura perturbacyjna $\alpha_5$}

Z~rozwini\k{e}cia perturbacyjnego, wsp\'o\l{}czynnik $\alpha_5$
wyra\.za si\k{e} przez amplitudy ogon\'ow $A_n, B_n$ ($u_n = r\,f_n \sim B_n\sin r + A_n\cos r$):
\begin{equation}
  \alpha_5 = B_5 + \tfrac{A_3^2}{2} + A_2 A_4 - B_2 A_2 A_3 + \tfrac{A_2^2(B_2^2 - B_3)}{2} - \tfrac{A_2^4}{8},
\end{equation}
gdzie $B_5$ jest amplitud\k{a} sin-ogow\k{a} 5-tego rz\k{e}du perturbacji.

\medskip\noindent Znane amplitudy (25+ cyfr):
\begin{align}
  B_2 &= I_{\cos} = \tfrac{1}{2} - \tfrac{\ln 3}{8} \quad \text{(\textbf{udow.})}, \notag \\
  A_2 &= \tfrac{\pi}{8} \quad \text{(\textbf{udow.})}, \notag \\
  B_3 &= P_{\cos} = 0.012615939290\ldots \quad \text{(30 cyfr)}, \notag \\
  A_3 &= -P_{\sin} = -0.215712018451\ldots \quad \text{(\textbf{NOWE}, 25 cyfr)}, \notag
\end{align}
gdzie $P_{\sin} = K_s^{(I)} - 2 K_s^{(II)}$ z~\textbf{udowodnionym}
$K_s^{(I)} = \pi/12$ (r\'ownanie~\eqref{eq:T6-KsI}) i~$K_s^{(II)} = 0.023043685\ldots$
(sw\'apowa dekompozycja, 25 cyfr).
Pozosta\l{}e potrzebne: $A_4, B_5$ (z~r\'owna\'n f$_4$, f$_5$).

\subsubsection*{Wynik perturbacyjny: $\alpha_5 = 0.02751$ (\textbf{5 cyfr})}

Rozwi\k{a}zuj\k{a}c r\'ownania perturbacyjne f$_2$, f$_3$, f$_4$, f$_5$
jako uk\l{}ad sprz\k{e}\.zony (8 komponent\'ow, scipy DOP853 do $R_{\max} = 5000$)
i~ekstrahuj\k{a}c amplitudy ogon\'ow z~dopasowania $u_n = r f_n \sim B_n\sin r + A_n\cos r$:
\begin{align}
  B_4 &\approx 0.08807, & A_4 &\approx -0.10039, \notag \\
  B_5 &\approx 0.00675, & A_5 &\approx -0.10732. \notag
\end{align}
B\l{}\k{a}d ekstrakcji $\sim 10^{-6}$ (oceniony z~b\l{}\k{e}du $B_3$ i~$A_3$ w~tych samych
r\'ownaniach).

\medskip\noindent Podstawiaj\k{a}c do formu\l{}y perturbacyjnej:
\begin{center}
\begin{tabular}{lr}
\hline
Cz\l{}on & Warto\'s\'c \\
\hline
$B_5$                         & $+0.00675$ \\
$A_3^2/2$                     & $+0.02327$ \\
$A_2 A_4$                     & $-0.03942$ \\
$-B_2 A_2 A_3$                & $+0.03072$ \\
$A_2^2(B_2^2 - B_3)/2$        & $+0.00917$ \\
$-A_2^4/8$                    & $-0.00297$ \\
\hline
$\alpha_5$                    & $\mathbf{+0.02751}$ \\
\hline
\end{tabular}
\end{center}
Zgodno\'s\'c z~pomiarem z~pe\l{}nego ODE ($\alpha_5 \approx 0.0275$, Richardson
$r_{\max} = 2000, 3000$) do $\sim 10^{-5}$. \textbf{Formu\l{}a perturbacyjna jest zweryfikowana.}

\medskip\noindent Analogicznie: $\alpha_4 = B_4 + A_2 A_3 - B_2 A_2^2/2 \approx -0.02460$
(zgodne z~pomiarem ODE $\sim -0.0246$).

\subsubsection*{Status}

Formu\l{}a perturbacyjna na $\alpha_5$ jest \textbf{zweryfikowana} ($\sim 5$ cyfr znacz\k{a}cych).
Dalsze zwi\k{e}kszenie precyzji wymaga analitycznej ewaluacji $B_4, A_4, B_5$
poprzez sw\'apow\k{a} dekompozycj\k{e} ca\l{}ek \'zr\'od\l{}owych
(analogicznie do $K_c^{(II)}$ i~$K_s^{(II)}$ --- mo\.zliwe, lecz z\l{}o\.zone).

\subsubsection*{Cz\k{e}\'sciowa dekompozycja analityczna \'zr\'od\l{}a $S_4$}

\'Zr\'od\l{}o $S_4$ rozk\l{}ada si\k{e} na pi\k{e}\'c cz\l{}on\'ow, z~kt\'orych
$T_1 = -f_1^2 (f_1')^2$ jest funkcj\k{a} \emph{wy\l{}\k{a}cznie} $f_1 = \sin(r)/r$.
Projekcje sinusowa i~kosinusowa tego cz\l{}onu s\k{a} obliczalne analitycznie:
\begin{align}
  \int_0^\infty \sin(t)\, t\, f_1^2\, (f_1')^2\, dt &= \tfrac{5\pi}{128} \quad \text{(\textbf{udow.})}, \\
  \int_0^\infty \sin(t)\, t\, f_1^3\, (f_1')^2\, dt &= \tfrac{\pi}{40} \quad \text{(\textbf{udow.})},
\end{align}
co wynika z~redukcji do standardowych ca\l{}ek sinc$^n$: $\int \operatorname{sinc}^4 = \pi/3$,
$\int \operatorname{sinc}^6 = 11\pi/40$.
Weryfikacja: PSLQ z~baz\k{a} zawieraj\k{a}c\k{a} $\pi$ daje dok\l{}adnie $40 A_5^{(f_1)} = \pi$
(40 cyfr). Projekcje kosinusowe (tj.~$B_4^{(T_1)}$, $B_5^{(f_1)}$) nie maj\k{a}
prostej formy zamkni\k{e}tej (PSLQ, 15 sta\l{}ych, do $10^6$).

% -------------------------------------------------------------------
\subsection{Podsumowanie i~otwarte kierunki}
\label{sec:T6-summary}
% -------------------------------------------------------------------

\begin{itemize}
  \item \textbf{Wynik g\l{}\'owny}: $\alpha_3 = \pi^2/128 + P_{\cos} = 0.089722223674\ldots$
        (30 cyfr), druga sta\l{}a TGP w~rozwini\k{e}ciu perturbacyjnym $\eta(\delta)$.
        \textbf{Hipoteza $\alpha_3 = \pi^2/110$ obalona}
        (r\'o\.znica $-1.45 \times 10^{-6}$, stabilna przy dps=25/40/60).
  \item \textbf{Metoda}: Sk\l{}adanie O($\delta^2$) + O($\delta^3$) perturbacji Besselowskich
        z~niezale\.zn\k{a} weryfikacj\k{a} numeryczn\k{a}: (a)~sw\'apowa dekompozycja
        z~kwadrosc\k{a} mpmath (dps=60), (b)~test zbie\.zno\'sci wieloprecyzyjny,
        (c)~sumowanie p\'o\l{}-okres\'ow z~akceleracj\k{a} Wynna-Epsilona.
  \item \textbf{Rozk\l{}ad}: $\alpha_3 = \pi^2/128 + P_{\cos}$ gdzie
        $\pi^2/128$ pochodzi z~O($\delta^2$) (\textbf{udowodnione}), a~$P_{\cos}$
        z~O($\delta^3$) (obliczone numerycznie do 30 cyfr, brak formy zamkni\k{e}tej).
  \item \textbf{Status analityczny}: $c_1 = 1 - \ln 3/4$ --- \textbf{udow.};
        $J_c = \ln(3)/8 - 1/2$, $J_s = \pi/8$ --- \textbf{udow.} (IBP);
        $K_c^{(I)} = (\ln 2 - 1)/6$ --- \textbf{udow.};
        $K_s^{(I)} = \pi/12$ --- \textbf{udow.} (sekcja~\ref{sec:T6-IBP});
        sw\'apowa dekompozycja $K_c^{(II)} = -A_1 - A_2 + A_3 - A_4$ z~jawnymi
        \emph{tail functions} $\Phi_i$ --- \textbf{udow.}
        (sekcja~\ref{sec:T6-swap}). $P_{\cos}$ zweryfikowane do 30 cyfr.
        PSLQ (baza 21 sta\l{}ych, max\_coeff $10^{10}$) nie~znajduje relacji.
        \textbf{Aktualizacja P4 (2026-04-19)}: rozszerzone poszukiwanie PSLQ
        z~baz\k{a} 24 sta\l{}ych przy max\_coeff do $10^{14}$ (skrypt
        \texttt{research/particle\_sector\_closure/ps1\_alpha3\_hunting.py},
        mpmath dps=40) tak\.ze nie~znajduje relacji liniowej. Na~tej podstawie
        $P_{\cos}$ zosta\l{}o promowane do nazwanej sta\l{}ej TGP\_Pcos
        (zob. Dodatek~\ref{app:P4-alpha3}).
  \item \textbf{Wsp\'o\l{}czynnik $\alpha_5 = 0.02751$}: Zweryfikowany perturbacyjnie
        (5~cyfr znacz\k{a}cych) przez rozwi\k{a}zanie uk\l{}adu sprz\k{e}\.zonego
        f$_2$--f$_5$ do $R_{\max} = 5000$. Formu\l{}a $\alpha_5 = B_5 + \frac{A_3^2}{2}
        + A_2 A_4 - B_2 A_2 A_3 + \frac{A_2^2(B_2^2-B_3)}{2} - \frac{A_2^4}{8}$
        jest \textbf{potwierdzona} zgodno\'sci\k{a} z~pomiarem ODE ($\sim 10^{-5}$).
        Analogicznie $\alpha_4 \approx -0.02460$. Identyfikacja PSLQ wymaga\l{}aby
        15+ cyfr, co wymaga sw\'apowej dekompozycji ca\l{}ek \'zr\'od\l{}owych.
  \item \textbf{Pr\'oba precyzji 10+ cyfr (2026-04-17)}: Przetestowano 7~podej\'s\'c
        (mpmath ODE dps=25/30, ca\l{}ki Greena on-the-fly, multiterm fit 8-16t,
        hybrid float64+mpmath). \textbf{Bariera 5--7 cyfr}: tail fitting ogranicza
        O($1/r^3$) residuum; Wynn-$\varepsilon$ nie konwerguje dla multi-frequency
        integrands ($f_2/f_3$-zale\.zne). Nowe wyniki analityczne:
        $\int \sin(t)\,t\,f_1^2(f_1')^2\,dt = 5\pi/128$,
        $\int \sin(t)\,t\,f_1^3(f_1')^2\,dt = \pi/40$ (PSLQ, 40 cyfr).
        Rzutowania kosinusowe tych ca\l{}ek NIE maj\k{a} prostej formy zamkni\k{e}tej.
  \item \textbf{Nast\k{e}pne kroki}: sw\'apowa dekompozycja Fubiniego dla~$B_4, A_4$
        (analogiczna do $K_c^{(II)}, K_s^{(II)}$) --- \textbf{jedyna \'scie\.zka}
        do 15+ cyfr niezb\k{e}dnych dla PSLQ.
\end{itemize}

\medskip
\textbf{Skrypty i~dane}:
\begin{itemize}
\item \texttt{research/brannen\_sqrt2/r6\_c2\_perturbative\_coeffs.py} --- pomiar $\alpha_2, \alpha_3$
\item \texttt{research/brannen\_sqrt2/r6\_c2\_ultra\_precision.py} --- ultra-prec do $r_{\max}=1600$
\item \texttt{research/brannen\_sqrt2/r6\_c3\_perturbative\_h\_and\_p.py} --- niezale\.zna weryfikacja $P_{\cos}$
\item \texttt{research/brannen\_sqrt2/r6\_c4\_proof\_Jc\_Js.py} --- dow\'od IBP $J_c, J_s$
\item \texttt{research/brannen\_sqrt2/r6\_c5\_KcI\_analytical.py} --- dow\'od $K_c^{(I)} = (\ln 2 - 1)/6$
\item \texttt{research/brannen\_sqrt2/r6\_c6\_KcII\_phi\_swap.py} --- sw\'ap dekompozycja $K_c^{(II)}$ z~$\Phi_i$
\item \texttt{research/brannen\_sqrt2/r6\_c7\_KcII\_high\_precision.py} --- dps=50 weryfikacja
\item \texttt{research/brannen\_sqrt2/r6\_c8\_moment\_integrals.py} --- tabela $M_1, M_2, M_3$
\item \texttt{research/brannen\_sqrt2/r6\_c9\_M2\_closed\_form.py} --- dow\'od $M_2, M_3$
\item \texttt{research/brannen\_sqrt2/r6\_c11\_alpha3\_ultra\_precision.py} --- $\alpha_3$ dps=50
\item \texttt{research/brannen\_sqrt2/r6\_c14\_alpha3\_precision\_convergence.py} --- test zbie\.zno\'sci dps=25/40/60 (\textbf{decyduj\k{a}cy})
\item \texttt{research/brannen\_sqrt2/r6\_c16\_quadosc\_error\_diagnostic.py} --- weryfikacja Wynn-Epsilon
\item \texttt{research/brannen\_sqrt2/r6\_c17\_pslq\_comprehensive.py} --- PSLQ 15-elementowa baza
\item \texttt{research/brannen\_sqrt2/r6\_c18\_pslq\_extended.py} --- PSLQ 21-elementowa baza (Cl$_2$, Ti$_2$, Li$_3$)
\item \texttt{research/brannen\_sqrt2/r6\_c20\_alpha5\_revised.py} --- pomiar $\alpha_5$ z~pe\l{}nego ODE (wi\k{e}ksze $\delta$)
\item \texttt{research/brannen\_sqrt2/r6\_c21\_alpha5\_richardson\_pslq.py} --- Richardson + PSLQ $\alpha_5$
\item \texttt{research/brannen\_sqrt2/r6\_c27\_KsI\_proof.py} --- dow\'od $K_s^{(I)} = \pi/12$ (40 cyfr)
\item \texttt{research/brannen\_sqrt2/r6\_c28\_KsII\_swap.py} --- $K_s^{(II)}$ sw\'ap $\to$ $P_{\sin}$, $A_3$ (25 cyfr)
\item \texttt{research/brannen\_sqrt2/r6\_c29\_coefficient\_summary.py} --- podsumowanie + PSLQ na nowych sta\l{}ych
\item \texttt{research/brannen\_sqrt2/r6\_c30\_A4B4\_perturbation\_ode.py} --- $A_4, B_4$ z~ODE perturbacyjnego ($R_{\max}=3000$)
\item \texttt{research/brannen\_sqrt2/r6\_c31\_B5\_fifth\_order.py} --- $B_5, A_5$ z~5-tego rz\k{e}du perturbacji
\item \texttt{research/brannen\_sqrt2/r6\_c33\_alpha5\_precise.py} --- $\alpha_5 = 0.02751$ (5 cyfr, $R_{\max} = 5000$)
\item \texttt{research/brannen\_sqrt2/r6\_c35\_B4\_decomposition.py} --- $\int\!\sin\!\cdot\!t\!\cdot\!f_1^2(f_1')^2 = 5\pi/128$ (PSLQ)
\item \texttt{research/brannen\_sqrt2/r6\_c38\_eta\_physical\_check.py} --- weryfikacja fizyczna $r_{21} = 206.54$
\item \texttt{research/brannen\_sqrt2/r6\_c46\_hybrid\_precision.py} --- hybrid float64+mpmath $\to$ 5--7 cyfr ($R=10000$)
\end{itemize}
